% 原书第3--8页至题型一结束；PDF第11--16页。人工对照原页转写。
\chapter{矩阵与线性变换}
高等代数内容包括：线性代数和多项式。本章重点关注矩阵与线性变换。
由于线性变换可用矩阵表示，本质上讲研究线性变换也就是研究矩阵。
因此我们讨论的主要内容包括矩阵的运算、线性变换的基本性质，以及矩阵的行列式、矩阵的特征值和特征向量、矩阵的秩、矩阵的迹等。
这些是线性代数的基础部分，特别地，矩阵与行列式是两个最基本的概念，熟练掌握它们的性质是最基本的要求。

\section{竞赛要点与难点}
\begin{enumerate}[label=（\arabic*）]
\item 矩阵的概念，矩阵的运算（加法、数乘、乘法、转置等运算）及其运算律。
\item 矩阵乘积的行列式、矩阵乘积的秩与其因子的秩的关系。
\item 矩阵的逆、伴随矩阵、矩阵可逆的条件，初等矩阵、初等变换、矩阵的等价标准形及逆矩阵的求解等。
\item 分块矩阵：分块矩阵及其运算与性质，分块初等矩阵、分块初等变换及其应用等。
\item 行列式的定义：2阶、3阶行列式的定义及对角线法则，排列与逆序数，$n$级行列式的定义等。
\item 行列式的性质。
\item 行列式的计算：根据定义计算行列式，利用行列式性质计算行列式，行列式行（列）展开、Laplace（拉普拉斯）展开定理等。
\item Cramer（克拉默）法则。
\item 矩阵的行秩、列秩、秩、矩阵的秩与其子式的关系。
\item 线性变换的定义、线性变换的运算、线性变换的矩阵。
\item 特征值与特征向量：线性变换和矩阵的特征值与特征向量、可对角化线性变换和矩阵的性质与判定、相似矩阵、相似不变量、Hamilton--Cayley（哈密顿--凯莱）定理等。
\item $\lambda$-矩阵：$\lambda$-矩阵、$\lambda$-矩阵在初等变换下的标准形、行列式因子、不变因子、初等因子，矩阵相似的条件等。
\item Jordan（若尔当）矩阵、Jordan标准形、有理标准形。
\end{enumerate}

\section{范例解析与精讲}
\subsection*{题型一\quad 矩阵的运算}
矩阵的运算中需要注意矩阵乘法不满足交换律，同时矩阵中有零因子存在，即两个非零矩阵乘积可以等于零矩阵。概括起来，以下几点需要特别注意：
\begin{enumerate}[label=（\arabic*）]
\item 一般地，二项式定理
\[
(A+B)^n=\sum_{i=0}^n\binom ni A^iB^{n-i}
\]
不一定成立。但若$A,B$可交换，即$AB=BA$，则上式成立。
\item 矩阵乘法没有消去律，即$CX=CY$（或$XC=YC$）推不出$X=Y$。但是有如下特别结论：
\begin{enumerate}[label=（\roman*）]
\item 一个矩阵$C$左可消当且仅当$C$为列满秩，即以下两条陈述等价：
\begin{enumerate}[label=（\alph*）]
\item 由$CX=CY$可推出$X=Y$。
\item $C$是列满秩，即$C$的列向量组线性无关或者$\rank(C)=C$的列数。
\end{enumerate}
\item 一个矩阵$C$右可消当且仅当$C$为行满秩，即以下两条陈述等价：
\begin{enumerate}[label=（\alph*）]
\item 由$XC=YC$可推出$X=Y$。
\item $C$是行满秩，即$C$的行向量组线性无关或者$\rank(C)=C$的行数。
\end{enumerate}
\item $A^2=E$推不出$A=E$或$A=-E$，即$A^2-E=(A+E)(A-E)=O$推不出$A+E=O$或$A-E=O$。
\end{enumerate}
\end{enumerate}

\ex{1.2.1}设$A,B$均是$n$阶方阵且满足$A^2=A$，$B^2=B$，$(A+B)^2=A+B$，求$AB$。
\sol 由$(A+B)^2=A+B$得到
\[A^2+AB+BA+B^2=A+B.\]
又$A^2=A$，$B^2=B$，所以
\[AB=-BA,\]
从而
\[A\cdot AB=-A\cdot BA,\]
即
\[AB=-ABA;\]
\[AB\cdot A=-BA\cdot A,\]
即
\[ABA=-BA\quad\text{或}\quad BA=-ABA,\]
因此
\[AB=BA.\]
由此得到$2AB=O$，故有$AB=O$。

\ex{1.2.3}设
\[
A=\begin{pmatrix}
a_1b_1&a_1b_2&\cdots&a_1b_n\\
a_2b_1&a_2b_2&\cdots&a_2b_n\\
\vdots&\vdots&&\vdots\\
a_nb_1&a_nb_2&\cdots&a_nb_n
\end{pmatrix},
\]
$k$为正整数，求$A^k$。
\sol $A$可分解为
\[
A=\begin{pmatrix}a_1\\a_2\\\vdots\\a_n\end{pmatrix}
\begin{pmatrix}b_1&b_2&\cdots&b_n\end{pmatrix}
=\alpha\beta^{\mathrm T},
\]
其中
\[
\alpha=\begin{pmatrix}a_1\\a_2\\\vdots\\a_n\end{pmatrix},\qquad
\beta=\begin{pmatrix}b_1\\b_2\\\vdots\\b_n\end{pmatrix},
\]
所以
\begin{align*}
A^k&=\underbrace{(\alpha\beta^{\mathrm T})(\alpha\beta^{\mathrm T})\cdots(\alpha\beta^{\mathrm T})}_{k\text{个}}\\
&=\alpha\underbrace{(\beta^{\mathrm T}\alpha)(\beta^{\mathrm T}\alpha)\cdots(\beta^{\mathrm T}\alpha)}_{k-1\text{个}}\beta^{\mathrm T}\\
&=\alpha(\beta^{\mathrm T}\alpha)^{k-1}\beta^{\mathrm T},
\end{align*}
而
\[\beta^{\mathrm T}\alpha=a_1b_1+a_2b_2+\cdots+a_nb_n,\]
故有
\begin{align*}
A^k&=(a_1b_1+a_2b_2+\cdots+a_nb_n)^{k-1}A\\
&=(a_1b_1+a_2b_2+\cdots+a_nb_n)^{k-1}
\begin{pmatrix}
a_1b_1&a_1b_2&\cdots&a_1b_n\\
a_2b_1&a_2b_2&\cdots&a_2b_n\\
\vdots&\vdots&&\vdots\\
a_nb_1&a_nb_2&\cdots&a_nb_n
\end{pmatrix}.
\end{align*}

\ex{1.2.5}设
\[
A=\begin{pmatrix}
2&0&0&0&0&0\\
3&2&0&0&0&0\\
0&3&2&0&0&0\\
0&0&0&1&2&0\\
0&0&0&0&1&2\\
0&0&0&0&0&1
\end{pmatrix},
\]
求$A^{100}$。
\sol 设
\[
A_1=\begin{pmatrix}2&0&0\\3&2&0\\0&3&2\end{pmatrix},\qquad
A_2=\begin{pmatrix}1&2&0\\0&1&2\\0&0&1\end{pmatrix}.
\]
则$A$可写为如下分块形式：
\[A=\begin{pmatrix}A_1&\\&A_2\end{pmatrix},\]
于是
\[A^{100}=\begin{pmatrix}A_1^{100}&\\&A_2^{100}\end{pmatrix}.\]
将$A_1,A_2$分解为
\[
A_1=2E+3S,\qquad\text{其中}\quad
S=\begin{pmatrix}0&0&0\\1&0&0\\0&1&0\end{pmatrix};
\]
\[
A_2=E+2H,\qquad\text{其中}\quad
H=\begin{pmatrix}0&1&0\\0&0&1\\0&0&0\end{pmatrix}.
\]
注意到$S$与$E$、$H$与$E$均可交换且
\[
S^2=\begin{pmatrix}0&0&0\\0&0&0\\1&0&0\end{pmatrix},\qquad
S^i=O\quad(i\ge3),
\]
\[
H^2=\begin{pmatrix}0&0&1\\0&0&0\\0&0&0\end{pmatrix},\qquad
H^i=O\quad(i\ge3).
\]
由二项式定理
\begin{align*}
A_1^{100}
&=(2E+3S)^{100}=\sum_{i=0}^{100}\binom{100}{i}(2E)^{100-i}(3S)^i\\
&=\binom{100}{0}(2E)^{100}
 +\binom{100}{1}(2E)^{99}\cdot3S
 +\binom{100}{2}(2E)^{98}(3S)^2\\
&=2^{100}E+300S\cdot2^{99}E
 +\frac{100\times99}{2}\cdot9S^2\cdot2^{98}E\\
&=2^{100}E+300\cdot2^{99}\cdot S+50\cdot99\cdot9\cdot2^{98}\cdot S^2\\
&=\begin{pmatrix}
2^{100}&0&0\\
3\cdot5^2\cdot2^{101}&2^{100}&0\\
11\cdot5^2\cdot3^4\cdot2^{99}&3\cdot5^2\cdot2^{101}&2^{100}
\end{pmatrix},
\end{align*}
\begin{align*}
A_2^{100}
&=(E+2H)^{100}
 =E^{100}+100\cdot E^{99}\cdot(2H)
 +\frac{100\times99}{2}\cdot E^{98}\cdot(2H)^2\\
&=E+200H+19800H^2\\
&=\begin{pmatrix}1&200&19800\\0&1&200\\0&0&1\end{pmatrix}.
\end{align*}
因此
\[
\setlength{\arraycolsep}{3pt}
A^{100}=\begin{pmatrix}
2^{100}&0&0&0&0&0\\
3\cdot5^2\cdot2^{101}&2^{100}&0&0&0&0\\
11\cdot5^2\cdot3^4\cdot2^{99}&3\cdot5^2\cdot2^{101}&2^{100}&0&0&0\\
0&0&0&1&200&19800\\
0&0&0&0&1&200\\
0&0&0&0&0&1
\end{pmatrix}.
\]

\par\medskip\noindent\textbf{总结}\quad 计算方阵的方幂一般有三种方法：
\begin{enumerate}[label=（\arabic*）]
\item 基于相似变换的计算法。更常见的情形是利用相似标准形（比如Jordan标准形）来进行计算。
\item 利用递推公式法或者说数学归纳法，如例1.2.3。
\item 基于方阵分解，然后利用二项式定理来进行计算，如例1.2.5。
\end{enumerate}
